% !PRINTED 260
\weberpage{260}

% Running head "Seventh Section." excluded as page furniture.

%<<CONT
If $\pi, \nu, \varrho$ are the number of positive, negative, vanishing squares into which $\varphi$ can be decomposed, then $\pi, \nu$ are the
numbers of positive and negative squares into which the
form $\Phi$ determined by (7) can be decomposed, and the determination of the numbers $\pi, \nu$ therefore need only be carried out for the latter function, whose determinant is different from
zero.

\begin{center}
\S.\,82.

\gesperrt{Quadratic forms with non-vanishing\\
determinant.}
\end{center}

In the investigation of the forms $\varphi(x_1, x_2 \ldots x_m)$ with non-vanishing determinant we make use of a theorem on determinants which we proved at the end of \S.\,28, and which
we wish to recall here.

If $A$ is a determinant of $m$ rows, and if $A_i^{(k)}$ are its
first, $A_{i,i'}^{k,k'}$ its second subdeterminants, then
\weq{1}{A_1^{(1)} A_i^{(i)} - A_1^{(i)} A_i^{(1)} = A A_{1,i}^{1,i}.}
If $A$ is a \gesperrt{symmetric determinant} different from zero, then $A_1^{(i)} = A_i^{(1)}$ and we can write formula (1)
thus:
\weq{2}{A_1^{(1)} A_i^{(i)} - A_1^{(i)2} = A A_{1,i}^{1,i}.}
In this, $i$ may denote any of the indices $2, 3 \ldots m$. If
$A_1^{(1)} = 0$ and $A_{1,i}^{1,i} = 0$, it follows from this that $A_1^{(i)} = 0$
must also hold, and from this we conclude that
\[
A_1^{(1)}, \; A_{1,2}^{1,2}, \; A_{1,3}^{1,3} \ldots A_{1,m}^{1,m}
\]
cannot vanish simultaneously, since otherwise
\[
A_1^{(2)}, \; A_1^{(3)} \ldots A_1^{(m)},
\]
and hence also, contrary to the hypothesis, $A$ would vanish (\S.\,22).

If we apply this to the determinant $R = R_m$ of our function $\varphi$, now assumed to be different from zero, it
follows that though the first principal subdeterminant $R_{m-1}$ may vanish,
not all the second principal subdeterminants $R_{m-2}$ can
vanish. The same mode of reasoning can be applied
when we replace $R$ by a non-vanishing $R_{m-1}, R_{m-2} \ldots$,
and we thus arrive at the following important theorem:
%CONT>>
